\section{Related work}\label{sec:related}

\paragraph{Stream clustering} The two-phase paradigm, online micro-clusters followed by offline macro-clustering, was introduced by CluStream~\cite{aggarwal2003clustream}, building on BIRCH's cluster features~\cite{zhang1996birch}. DenStream~\cite{cao2006denstream} added exponential decay and density-based macro-clustering, and DBSTREAM~\cite{hahsler2016dbstream} a shared-density graph; see~\cite{zubaroglu2021review} for a survey. FedCAST uses decayed CFs as its unit of exchange because they are additive, cheap to send, and give the k-means cost in closed form, which is what makes an objective-linked trigger computable on the device.

\paragraph{Distributed stream clustering} Cormode et al.~\cite{cormode2007divide} maintain a continuous k-center clustering over distributed streams with approximation guarantees. DGClust~\cite{gama2011dgclust} sends grid-cell state changes from sensors, and Tran and Sattler~\cite{tran2013communication} send micro-clusters from remote sites to a coordinator. These works reduce communication with fixed, objective-agnostic rules (state change, period or threshold). They have no byte budget, do not adapt to link conditions, do not treat statistical heterogeneity, and are evaluated in simulation without a messaging substrate. Coreset-based methods~\cite{balcan2013distributed} offer strong guarantees for static data but no cheap incremental delta.

\paragraph{Federated clustering} k-FED~\cite{dennis2021kfed} performs one-shot federated k-means and shows that heterogeneity \emph{helps} when each device holds few clusters. FFCM~\cite{stallmann2022ffcm} federates fuzzy c-means, and AFCL~\cite{zhang2025afcl} handles asynchrony and an unknown number of clusters. All of them assume static local datasets. We run k-FED periodically as a strong baseline and report the regimes in which it wins.

\paragraph{Federated learning on streams} FedCluLearn~\cite{angelova2025fedclulearn} uses micro-cluster indexing to fight forgetting in federated continual \emph{supervised} learning. Silva et al.~\cite{silva2022federated} and DFAS~\cite{li2025dfas} address federated anomaly detection on streams. None of these targets clustering of the data itself under an explicit communication budget.

\paragraph{Event-triggered and communication-efficient learning} Communication efficiency is central in federated learning~\cite{mcmahan2017fedavg,kairouz2021advances}. Event-triggered schemes send when the norm of a parameter change exceeds a threshold, possibly adapted to device resources~\cite{ghosh2021eventgrad,zehtabi2022event}. FedCAST differs in three ways: its trigger is measured in the clustering objective and bounds the server's error (Propositions~\ref{prop:cost} and~\ref{prop:centroid}); it enforces a byte budget with a guarantee (Proposition~\ref{prop:budget}); and it prioritises \emph{which} parts of the model to send.

\paragraph{Kafka for federated ML} Kafka-ML~\cite{martin2022kafkaml} and its federated extension~\cite{carnerero2024kafkaml} use Kafka as the data and model transport for neural-network training. In FedCAST, Kafka's ordering, idempotence, replay and compaction semantics are part of the algorithm's recovery and exactly-once-state argument, and we measure them under crashes.
